\section{Introduction}

World modeling and video generation have emerged as central challenges in artificial intelligence, requiring the synthesis of temporally coherent and photorealistic sequences conditioned on semantically rich inputs such as natural language, static imagery, or short video clips. This task resides at the intersection of spatial perception and temporal reasoning, with profound implications for fields including robotics, embodied artificial intelligence, interactive media, and scientific simulation. As video becomes a dominant modality for both human communication and machine understanding, the demand for generative models that are not only high-fidelity and computationally efficient, but also causally consistent and compatible with streaming applications, has become increasingly urgent.

Building on the remarkable success of diffusion~\citep{sohl2015deep, ho2020denoising, song2020score} and flow-matching frameworks~\citep{lipman2022flow, liu2022flow} in image generation, recent research has increasingly focused on extending these approaches to video synthesis. However, most large-scale video diffusion models continue to rely on globally conditioned denoising architectures that process the entire temporal sequence simultaneously. These models typically employ uniform noise levels and require full-sequence access during inference. Such designs disregard the causal structure inherent to temporal data, rendering them suboptimal for scenarios requiring streaming, real-time interaction, or autoregressive generation.

To overcome these limitations, we present \magi: a large-scale diffusion-based generative model that produces video through the \textit{autoregressive} generation of temporally segmented chunks, each consisting of a fixed-length sequence of consecutive frames. This chunk-wise approach offers a principled trade-off between causal modeling and temporal abstraction, enabling the model to capture mid-range temporal dependencies while maintaining strict left-to-right temporal consistency. Training is conducted at the chunk level with temporally \textit{progressive} noise levels, resulting in a model that is both autoregressively structured and adaptable in its conditional generation capacity.

\magi adheres strictly to causal constraints and facilitates real-time, streaming-compatible video synthesis that approximates multi-step diffusion trajectories with reduced-step, chunk-level predictions. This is enabled by a Transformer~\citep{transformer} backbone specifically designed for bidirectional spatial and causal temporal denoising, supported by a carefully engineered training infrastructure. Central to this infrastructure is a novel distributed attention mechanism (MagiAttention) tailored for ultra-long autoregressive contexts, along with a scalable execution framework optimized for low-latency, parallelized inference. These core components are further augmented by a robust data curation pipeline that supports multi-stage training and dynamically adapts the data distribution based on ongoing model evaluation. Together, these architectural and algorithmic advances empower \magi to deliver efficient, scalable, and controllable video generation. Notably, the inference-time peak resource usage of \magi is independent of the total video length, as each chunk is processed with a fixed computational and memory footprint. This makes \magi particularly suitable for low-latency, memory-efficient applications. The largest variant of the model comprises 24 billion parameters and supports context lengths of up to 4 million tokens, demonstrating the scalability and robustness of the framework.

We evaluate \magi using both internal metrics and publicly available benchmarks, with a particular focus on the image-to-video (I2V) generation task. Our evaluation protocol assesses prompt fidelity, temporal coherence, and subject integrity. On VBench-I2V~\citep{huang2024vbenchcomprehensiveversatilebenchmark} and Physics-IQ Benchmark~\citep{physicsiqbenchmark}, \magi achieves substantial improvements over previous models, especially in its ability to synthesize complex motion, preserve semantic alignment, and model physically plausible interactions.

In summary, \magi establishes a scalable and autoregressive foundation for diffusion-based video synthesis. By integrating architectural innovations, high-throughput inference techniques, and a comprehensive data processing framework, \magi bridges the gap between high-quality generative performance and real-time applicability. The complete inference codebase and pre-trained models are publicly accessible at \href{https://github.com/sandai-org/magi-1}{\texttt{magi-source}}, the distributed attention available at \href{https://github.com/sandai-org/magiattention}{\texttt{magi-attention}}, and a live demonstration available at \href{https://sand.ai}{\texttt{magi-product}}.



